\section{Bounded relational characterization}
\label{sec:alloy}

\subsection{Model structure}

The Alloy model represents records, authoritative fields, linearly ordered versions, evidence content and locators, candidates, certificates and primitive identities, principals, authorizations, transitions, and states. A state maps each committed record/field pair to at most one value and one source transition. Once a record has an authoritative snapshot, both maps cover exactly its declared field domain.

Three event families define the bounded histories. A \code{MachineEvent} may extend candidate, certificate, and evidence sets but frames authoritative values, sources, versions, transitions, and authorizations, encoding P0. A \code{BatchAdmissionEvent} contains a nonempty item set sharing an event envelope. A \code{TamperEvent} freezes values and source pointers while changing only the transition set, giving P6 a deliberately narrow corruption boundary.

The base effect advances one record version and applies item values and sources while intentionally leaving transition structure underconstrained. Full supplies unique item fields; initial or later complete-snapshot shape; preexisting certificate and evidence; record, field, evidence-identity, evidence-owner, certificate-version, freshness, authorization-certificate, principal, and authorization-value bindings; and an exact item--transition bijection. Keeping these conjuncts separate permits one encoded distinction to be removed without silently changing the others.

\subsection{From information classes to Alloy checks}

Table~\ref{tab:distinction-evidence} connects the representation-independent classes of \cref{eq:distinction-basis} to the concrete relational operators and executable evidence. Several low-level conjuncts may encode one high-level class; the eleven ablations are therefore not eleven separate novelty claims.

\begin{table}[t]
\centering
\footnotesize
\caption{Mapping from failure-distinguishing information classes to bounded and concrete evidence.}
\label{tab:distinction-evidence}
\begin{tabularx}{\textwidth}{@{}l>{\raggedright\arraybackslash}p{0.28\textwidth}>{\raggedright\arraybackslash}p{0.29\textwidth}>{\raggedright\arraybackslash}X@{}}
\toprule
Class & Safe/unsafe witness & Alloy support & Concrete support \\
\midrule
$D_C$ & exact candidate / equal-valued contextual substitute & field, evidence-identity, evidence-owner, record, authorization-certificate, certificate-preexistence, and evidence-preexistence ablations & record, field, evidence, and certificate substitution cases \\
$D_V$ & retained Correction / wrong attribution of the same final value & legal Correction witnesses and authorization-value ablation & singleton, all-field, and mixed Correction lifecycles with immutable candidate checks \\
$D_F$ & current / superseded authorization & certificate-version and freshness ablations; stale-rejection assertions & stale replay, contending confirmation, and compare-and-swap rejection \\
$D_B$ & atomic / partial or fragmented successor & transition-bijection ablation and whole-batch framing assertions & exact changed-field validation, failpoints, rollback, and concurrency cases \\
$D_S$ & complete / missing, replaced, or duplicate source & source-anchor attack witnesses and P6 trace assertions & source-corruption catalogue and reverse-trace validation \\
\bottomrule
\end{tabularx}
\end{table}

Each ablation pair contains at least two items: one violates only the selected conjunct and another remains Full-valid. The same attempted item set is admitted by the reduced relation and rejected with authoritative-state stutter by Full. This tests the sensitivity of the selected relational conjunct. The observation-level irredundancy argument is the separate paired-history construction in \cref{sec:contract}.

\subsection{Legal, unsafe, and preservation commands}

Six legal-witness commands cover multi-field Accept, Correction, mixed Accept/Correction, same-value admission, initial snapshot, and singleton admission. Eleven paired ablations remove the conjuncts mapped in \cref{tab:distinction-evidence}. Principal membership remains fixed; the unused \code{ABL\_7} marker is excluded because host identity trust is an external assumption.

Three source-attack commands start from a legal Full prefix and then remove a source anchor, replace it with another transition atom, or add a duplicate source-version transition; each requires the trace to become incomplete. Thirteen assertions per scope cover invariant preservation, P0/P1/P3/P4/P5 regressions, whole-batch effect and rejection framing, bidirectional singleton correspondence, Full rejection of an ablated attempt, and P6 incompleteness.

Nine assertions guard definitions and their immediate consequences: P0/P1/P3/P4/P5, whole-batch effects, rejection framing, and the two P6 checks. The other four check state-invariant preservation, the two singleton correspondences, and Full rejection of an ablated attempt. This grouping distinguishes roles in the validation argument; thirteen UNSAT commands are not thirteen independent discoveries or proofs. The command set tests legal and attack reachability, preservation, and sensitivity to individual conjunct removal within each scope.

\subsection{Scopes and interpretation}

The S1 profile bounds principal structural domains at two records, three fields, six versions, up to six admission items/transitions for most commands, and four states/events. S2 increases representative bounds to three records, four fields, eight versions, up to eight admission items/transitions, and five states/events. Individual commands adjust Value, Evidence, or Transition bounds when a witness needs an extra atom. The artifact contains the executable scopes; this summary is not a replacement for them.

% Generated from the frozen r27 Alloy receipts.
% Source: formal/alloy/batch/command_manifest.tsv and formal/alloy/batch/raw-results/2026-09-10-trace-positive/command_results.json
% r27 extension: +1 positive trace-completeness assertion per profile (34 commands per profile; 68 total).
\begin{table}[t]
\centering
\caption{Bounded Alloy outcomes in the two declared scopes. SAT denotes a found witness; UNSAT denotes no instance within the encoded scope.}
\label{tab:alloy-outcomes}
\begin{tabularx}{\linewidth}{@{}l*{5}{>{\centering\arraybackslash}X}>{\centering\arraybackslash}p{0.15\linewidth}@{}}
\toprule
Profile & Legal SAT & Ablation SAT & Attack SAT & Total SAT & Total UNSAT & Commands \\
\midrule
S1 & 6 & 11 & 3 & 20 & 14 & 34 \\
S2 & 6 & 11 & 3 & 20 & 14 & 34 \\
\bottomrule
\end{tabularx}
\end{table}


An UNSAT assertion means that the Analyzer found no counterexample within the encoded command scope. A SAT run means that it found the requested witness \citep{jackson2002alloy,torlak2007kodkod}. Singleton contract/effect checks additionally guard the batch lift against silently changing the historical one-item semantics.
